\section{Related and Future Work}
%\section{Conclusions and future work}
\label{s:cf}
              

\dsOLD{maybe this section could be merged with the conclusion ``Related  and future work'';
  also we could try to have a stronger structure (maybe subsections, not sure)}

Analogies between game semantics and $\pi$-calculus, as semantic frameworks in which
\emph{names} are central, have been pointed out from the very beginning of game
semantics. 
In the pioneering work~\cite{10.1145/224164.224189}, the authors 
obtain 
 a translation of PCF terms into the $\pi$-calculus from 
a game model of PCF by representing 
 \emph{strategies} 
(the denotation of PCF terms in the game model) as processes of the $\pi$-calculus.
The encoding bears similarities with Milner's, though they are syntactically rather
different (``it is clear that the two are conceptually quite unrelated'',
\cite{10.1145/224164.224189}).   The connection has been  developed in various papers,
e.g.,
\cite{DBLP:journals/entcs/Honda02,10.1016/S0304-3975-99-00039-0,10.5555/788020.788890,yoshida2020game,10.1145/3290340}.
% in certain cases (notably
% ) the definition of the game models has been inspired 
Milner's encodings into the $\pi$-calculus have sometimes  been a source of inspiration  in
the definition of the game semantics models (e.g.,  transporting 
 the work 
\cite{10.1145/224164.224189}, in call-by-name, onto call-by-value
 \cite{10.1016/S0304-3975-99-00039-0}).  
In~\cite{yoshida2020game}, a  typed variant of  $\pi$-calculus, influenced by 
differential linear logic~\cite{ehrhard2018introduction},
is introduced as a metalanguage to represent game models.
\iflong 
 The approach is illustrated on a
a higher-order call-by-value language with shared memory concurrency, whose game model is build by first translated it
to this metalanguage, then using the game model of this metalanguage.
\fi
In~\cite{10.1109/LICS.2015.20}, games are defined using algebraic operations on sets of
traces, and used to  prove type soundness of a simply-typed call-by-value
$\lambda$-calculus with effects.
% like $\callcc$, references and fork.  
 Although the calculus of traces employed is not a $\pi$-calculus 
(e.g., being defined from operators and relations over trace sets rather than from
syntactic process constructs), there are similarities, which would be interesting to
investigate. 



Usually in the above  papers
the source language is a form of $\lambda$-calculus, 
that is interpreted into game semantics,  
and the 
 $\pi$-calculus  (or dialects of it)  is used to
represent the resulting strategies and games. 
Another goal has been to shed light on typing
disciplines  for 
$\pi$-calculus processes, by transplanting conditions on strategies such as well-bracketing
and innocence into appropriate typings for the $\pi$-calculus (see, e.g., 
\cite{berger2001sequentiality,YBH11,HirschkoffPS21}). 


In contrast with the above works,  where analogies between game semantics and
$\pi$-calculus are used to better understand one of the two models 
(i.e., explaining  game semantics in terms process interactions,  or enhancing type systems for
processes following structures in game semantics), 
in the  present paper we have carried out a direct comparison between 
the two models (precisely OGS and \piI). 
For this we have started from
 the (arguably natural)  representations of the $\lambda$-calculus 
into OGS and \piI  (the latter being Milner's encodings).
Our goal was
understanding the  relation between  the behaviours of the terms in the
two models, and   transferring techniques and results between them.

Technically, our work  builds on
\cite{DurierHS18,Durier20}, where a  detailed analysis of the behaviour of Milner's
call-by-value encoding is carried out using proof techniques for
$\pi$-calculus based on unique-solution of equations.
Various results in \cite{DurierHS18,Durier20} are essential to our own
 (the observation that Milner's
encodings should be interpreted in \piI rather than the full $\pi$-calculus is also from
\cite{DurierHS18,Durier20}).

An OGS %LTS 
for the call-by-name $\lambda\mu$-calculus is introduced
in~\cite{DBLP:conf/fossacs/Laird20}. It 
 is different, once restricted to the fragment without the $\mu$-binder, to the one we
 consider in Section~\ref{s:cbnNEW}.  It might correspond  
 to a different way of encoding  call-by-name into the $\pi$-calculus
where $\lambda$-abstractions are interpreted as input-prefixed processes, as e.g., in 
 Milner's original encoding~\cite{Mil92s}.



%We conjecture that this LTS would rather correspond to Milner's original encoding of the lazy $\lambda$-calculus in the $\pi$-calculus~\cite{Mil92s},
%where $\lambda$-abstractions are interpreted directly as inputs.


Bisimulations over  OGS terms, and tensor products of configurations,
were introduced in~\cite{10.1145/2603088.2603150},
in order to provide a framework to study compositionality properties of OGS.
In our case the compositionality result of OGS is derived from the correspondence with the
$\pi$-calculus.
% \iflong
%  Their notion of tensor on states corresponds to our notion of composition of configurations.
%  They use an abstract notion of \emph{linker} to define \emph{transfer} between strategies.
%  The link $\fwd{a}{b}$ should fit this abstract definition of a linker.
% \fi
In~\cite{sakayori2019categorical}, a correspondence between an i/o typed asynchronous $\pi$-calculus 
and a computational $\lambda$-calculus with channel communication is established, 
using a common categorical model (a compact closed Freyd category).
It would be interesting to see if our concurrent operational game model could be equipped with this categorical semantics.

\emph{Normal form} (or \emph{open}) bisimulations~\cite{lassentrees,Stovring-Lassen},
as game
semantics, 
manipulate open terms, and sometimes make use of
environments or stacks of evaluation contexts (see e.g., the recent
work~\cite{biernacki2019complete}, where  a fully abstract  normal-form 
bisimulation for a $\lambda$-calculus with store is obtained).  

Notice that contextual equivalence for such a language corresponds to complete well-bracketed 
trace equivalence.
It would be interesting to see if, once the reasoning on diverging configurations removed,
the enf-bisimulations obtained there could be seen as up-to composition bisimulations
over the OGS LTS \cutLMCSsecondRound{for the $\lambda\rho$-calculus presented in Appendix~\ref{a:LambdaRho},
that is fully abstract for this language}.

There are also works that build
game models directly for the $\pi$-calculus.
In~\cite{laird2005game}, a game models for the simply-typed asynchronous $\pi$-calculus 
is introduced, and proved fully abstract for may-equivalence.
In~\cite{10.1007/11944836_38}, this construction was extended to a model of Asynchronous Concurrent ML.
These models rely on a sequential representation of interactions as traces.
More recently, causal representation of game models of the $\pi$-calculus have been proposed.
In~\cite{sakayori2017truly}, a truly concurrent game model for the asynchronous $\pi$-calculus is introduced,
using a directed acyclic graph structure to represent interactions.
%
A game model for a synchronous $\pi$-calculus is defined in~\cite{eberhart2015intensionally}, and shown to
be fully abstract for fair testing equivalence. 
A correspondence between a synchronous $\pi$-calculus with session types and 
concurrent game semantics~\cite{winskel2017games}
is given in~\cite{10.1145/3290340}, relating 
games (represented as arenas) to  session types, and 
strategies (defined as coincident event structures) to processes.


We have exploited the full abstraction results between OGS and \piI
 to transport a few  up-to techniques for
  bisimulation  from
  \piI onto OGS.  However, in \piI there are various other such techniques, even a theory of
 bisimulation enhancements.  We would like to see which other techniques could be useful
 in OGS,  possibly transporting  the theory of enhancements itself. 
We would like also to study the possible game-semantics counterpart of recent work on a
 representation of functions as processes \cite{SakayoriS23} in which different tree structures of the
$\lambda$-calculus (Lévy-Longo  trees,  Böhm trees,  Böhm trees with infinite $\eta$) are obtained
by modifying the semantics of the certain special processes used as a parameter of the
representation
 (and related to the \emph{link} processes of
Section~\ref{s:encoLpi} and \ref{s:cbnNEW}, called \emph{wires} in~\cite{SakayoriS23}). 